%============================================================
% CH2 · P3 · PROBLEM 1 — MILITARY PROBLEM (CF 1006E)
%============================================================
\subsection{Problem 2.3.1: Military Problem}
\problemlabel{2.3.1}{ssec:ch2-military}

\problemcard{2.3.1}{Military Problem}{Codeforces 1006E}{https://codeforces.com/problemset/problem/1006/E}{\exmedium\ \cf{1600}}{One pre-order of the whole tree; the command list of $u$ is the block \texttt{order[pos[u] .. pos[u]+sz[u]-1]}}

% local layout of the sample tree
\tikzset{
    g2p3m-v1/.style=gv, g2p3m-v2/.style=gv, g2p3m-v3/.style=gv, g2p3m-v4/.style=gv, g2p3m-v5/.style=gv,
    g2p3m-v6/.style=gv, g2p3m-v7/.style=gv, g2p3m-v8/.style=gv, g2p3m-v9/.style=gv,
    g2p3m-e/.style={gde},
    g2p3m-T/.pic={
        \node[g2p3m-v1] (m1) at (2.4,0) {1};
        \node[g2p3m-v2] (m2) at (0.6,-1.2) {2};
        \node[g2p3m-v3] (m3) at (2.4,-1.2) {3};
        \node[g2p3m-v4] (m4) at (4.2,-1.2) {4};
        \node[g2p3m-v5] (m5) at (1.6,-2.4) {5};
        \node[g2p3m-v7] (m7) at (3.2,-2.4) {7};
        \node[g2p3m-v6] (m6) at (1.0,-3.6) {6};
        \node[g2p3m-v8] (m8) at (2.2,-3.6) {8};
        \node[g2p3m-v9] (m9) at (3.2,-3.6) {9};
        \foreach \a/\b in {1/2, 1/3, 1/4, 3/5, 3/7, 5/6, 5/8, 7/9} \draw[g2p3m-e] (m\a) -- (m\b);
    }
}

\subsubsection*{1. Problem statement}

An army of $n$ officers forms a tree rooted at officer 1; officer $i \ge 2$ has the direct superior $p_i < i$. When officer $u$ issues a command it spreads by DFS: the current officer hands it to his direct subordinate with the \textbf{smallest index} that has not received it yet, waits until that subordinate has finished spreading it through his own subtree, then picks the next one; when no subordinate is left, he is done. Answer $q$ independent queries $(u_i, k_i)$: the $k_i$-th officer to receive the command issued by $u_i$ ($u_i$ himself is the first), or $-1$ if fewer than $k_i$ officers receive it.

\textbf{Constraints:} $2 \le n \le 2\cdot10^5$, $1 \le q \le 2\cdot10^5$, $1 \le p_i < i$, $1 \le u_i, k_i \le n$.

\begin{table}[H]
\centering\small
\begin{tabular}{@{}lll@{}}
\toprule
\textbf{Input} & \textbf{Official output} & \textbf{Our output} \\
\midrule
\texttt{9 6 / 1 1 1 3 5 3 5 7 / 3 1 / 1 5 / 3 4 / 7 3 / 1 8 / 1 9} & \texttt{3 6 8 -1 9 4} & \texttt{3 6 8 -1 9 4} \\
\bottomrule
\end{tabular}
\end{table}

\begin{diagrambox}
\begin{tikzpicture}[x=1cm, y=0.9cm]
    \begin{scope}
        \pic {g2p3m-T};
        \node[font=\scriptsize\bfseries] at (2.4,0.8) {the sample: $p = (1,1,1,3,5,3,5,7)$};
    \end{scope}
    \begin{scope}[xshift=5.5cm, g2p3m-v3/.style=gvsrc, g2p3m-v5/.style=gvseen, g2p3m-v6/.style=gvseen,
                  g2p3m-v8/.style=gvdst, g2p3m-v7/.style=gvseen, g2p3m-v9/.style=gvseen]
        \pic {g2p3m-T};
        \node[font=\scriptsize\bfseries] at (2.4,0.8) {command from 3: $3, 5, 6, \mathbf{8}, 7, 9$};
        \foreach \v/\k/\p in {3/1/right, 5/2/left, 6/3/below, 8/4/below, 7/5/right, 9/6/below}
            \node[gdist, text=accentorange, \p=0.6mm of m\v] {\tiny\k};
        \node[note, text width=3.3cm, align=left, anchor=west] at (4.9,-2.6) {query $(3, 4)$: the 4th officer is \textbf{8}; query $(7, 3)$: only $7, 9$ receive it, answer $-1$};
    \end{scope}
\end{tikzpicture}
\figcap{The sample tree and the spreading from officer 3 (small numbers: receiving order). Children are always taken by increasing index.}
\end{diagrambox}

\subsubsection*{2. The graph model}

\begin{modelbox}[Graph Model --- Military Problem]
\renewcommand{\arraystretch}{1.35}
\begin{tabularx}{\linewidth}{@{}>{\bfseries}l X@{}}
Vertices & the $n$ officers \\
Edges & tree edges $p_i \to i$, rooted at 1; children ordered by index \\
Weights & none \\
Question & the $k$-th vertex of the pre-order of the subtree of $u$ (with this child order) \\
Size & $n, q \le 2\cdot10^5$: \bigO{n} once plus \bigO{1} per query; the tree can be a path ($2\cdot10^5$ deep) \\
\end{tabularx}
\end{modelbox}

\begin{diagrambox}
\begin{tikzpicture}[x=1cm, y=1cm]
    \foreach \v/\i in {1/0, 2/1, 3/2, 5/3, 6/4, 8/5, 7/6, 9/7, 4/8}{
        \node[gcell, minimum size=8mm, font=\small\bfseries] (c\i) at (\i*0.95,0) {\v};
        \node[font=\tiny, text=darkgray] at (\i*0.95,-0.6) {pos \i};
    }
    \draw[primaryblue, line width=1.4pt, decorate, decoration={brace, amplitude=5pt}] (-0.45,0.5) -- (8.05,0.5)
        node[midway, above=5pt, font=\scriptsize, text=primaryblue] {command from 1: all 9 (sz 9)};
    \draw[accentorange, line width=1.3pt, decorate, decoration={brace, amplitude=4pt, mirror}] (1.45,-0.85) -- (7.1,-0.85)
        node[midway, below=4pt, font=\scriptsize, text=accentorange] {command from 3: pos $2..7$ (sz 6)};
    \draw[softpurple, line width=1.3pt, decorate, decoration={brace, amplitude=4pt}] (5.25,0.45) -- (7.1,0.45)
        node[midway, above=4pt, font=\scriptsize, text=softpurple] {from 7};
\end{tikzpicture}
\figcap{The pre-order of the whole tree, \texttt{order[]}. Every officer's command list is a contiguous block of it, starting at the officer himself.}
\end{diagrambox}

\subsubsection*{3. Key observations}

\begin{itemize}[leftmargin=*]
    \item \textbf{The spreading is a DFS pre-order} with children in increasing order. A command from $u$ reaches exactly the subtree of $u$, in the same relative order in which a DFS from the root would visit it --- the DFS from the root, once it enters $u$, does exactly what a DFS from $u$ does until it leaves $u$.
    \item \textbf{So all $n$ lists are blocks of one array.} By the pre-order invariant of the pattern, $u$'s list is $\texttt{order}[\mathrm{pos}[u]], \dots, \texttt{order}[\mathrm{pos}[u]+\mathrm{sz}[u]-1]$. The answer is $\texttt{order}[\mathrm{pos}[u]+k-1]$ if $k \le \mathrm{sz}[u]$, else $-1$.
    \item \textbf{Children come sorted for free.} Reading $p_2, p_3, \dots$ and appending $i$ to \texttt{ch[$p_i$]} puts the children in increasing order.
    \item \textbf{$p_i < i$ makes sizes trivial:} for $i = n$ down to $2$, every child is finished before its parent is read.
    \item \textbf{Depth.} $p_i = i - 1$ gives a path of $2\cdot10^5$ officers: the DFS must be iterative.
\end{itemize}

\subsubsection*{4. Approaches}

\begin{bruteforcebox}
\textbf{Idea: simulate.} For each query run the spreading from $u_i$ literally and stop after $k_i$ officers.

\begin{lstlisting}[style=cpp]
#include <bits/stdc++.h>
using namespace std;

// Brute force: simulate the spreading for every query, stop after k officers.
vector<vector<int>> ch;
vector<int> got;

void spread(int a, int k) {
    if ((int)got.size() == k) return;              // the k-th officer is known: stop early
    got.push_back(a);
    for (int b : ch[a]) spread(b, k);              // children are stored in increasing order
}

int main() {
    int n, q;
    cin >> n >> q;
    ch.assign(n + 1, vector<int>());
    for (int i = 2; i <= n; i++) {
        int p;
        cin >> p;
        ch[p].push_back(i);
    }
    while (q--) {
        int u, k;
        cin >> u >> k;
        got.clear();
        spread(u, k);
        cout << ((int)got.size() == k ? got[k - 1] : -1) << "\n";
    }
}
\end{lstlisting}

\begin{complexitybox}
\textbf{Time:} \bigO{n} per query, \bigO{nq} $= 4\cdot10^{10}$ in the worst case (a star with $u = 1$, $k = n$ on every query). \textbf{Space:} \bigO{n}, but the recursion is as deep as the tree --- on the path input it also crashes. \textbf{Verdict:} \tle{}; it is the reference for the random tests.
\end{complexitybox}
\end{bruteforcebox}

\begin{algorithmbox}[Approach 2: one pre-order, one block per officer]
\textbf{Idea.} Compute the pre-order of the whole tree once, iteratively: pop a vertex, give it the next position, push its children in \emph{reverse} order so that the smallest is popped next. Then sizes bottom-up, and each query is one array access.

\begin{lstlisting}[style=cpp]
#include <bits/stdc++.h>
using namespace std;

// Military Problem (CF 1006E): the command from u spreads in DFS pre-order of u's subtree
// (children by increasing index). One pre-order of the whole tree answers every query:
// u's list is the block order[pos[u] .. pos[u] + sz[u] - 1].
int main() {
    ios::sync_with_stdio(false);
    cin.tie(nullptr);
    int n, q;
    cin >> n >> q;
    vector<vector<int>> ch(n + 1);
    for (int i = 2; i <= n; i++) {
        int p;
        cin >> p;
        ch[p].push_back(i);                    // i increases: children lists come out sorted
    }
    vector<int> order, pos(n + 1), sz(n + 1, 1);
    vector<int> st(1, 1);
    while (!st.empty()) {                      // iterative pre-order (a path of 2e5 is legal)
        int u = st.back(); st.pop_back();
        pos[u] = order.size();
        order.push_back(u);
        for (int j = (int)ch[u].size() - 1; j >= 0; j--)
            st.push_back(ch[u][j]);            // reverse: the smallest child is popped first
    }
    for (int u = n; u >= 1; u--)               // p_i < i: every child is finished before u
        for (int c : ch[u]) sz[u] += sz[c];
    while (q--) {
        int u, k;
        cin >> u >> k;
        cout << (k <= sz[u] ? order[pos[u] + k - 1] : -1) << "\n";
    }
}
\end{lstlisting}

\begin{complexitybox}
\textbf{Time:} \bigO{n} for the pre-order and the sizes, \bigO{1} per query: \bigO{n + q}. \textbf{Space:} \bigO{n}, no recursion. \textbf{Verdict:} \ac{}-ready: on the path $p_i = i-1$ with $2\cdot10^5$ random queries the whole run, input and output included, takes under a second on this machine.
\end{complexitybox}
\end{algorithmbox}

\begin{invariantbox}[Why the block is the command list]
\textbf{Claim.} Popping from the stack produces a pre-order of the tree with children in increasing order, and the vertices popped from the moment $u$ is popped until the stack is back to the height it had below $u$ are exactly $u$'s subtree.

When $u$ is popped, its children are pushed above everything that was on the stack, smallest on top. Everything pushed later lies above them, so nothing below $u$'s children is popped before all of $u$'s children --- and, by induction, all their descendants --- have been popped. Hence $u$'s subtree is popped as one consecutive run that starts with $u$, and inside it the children's subtrees come one after another in increasing order of the child: the spreading order. The run has length $\mathrm{sz}[u]$, which the bottom-up loop computes ($p_i < i$, so when $u$ is processed all its children already hold final sizes). \hfill$\blacksquare$
\end{invariantbox}

\begin{alternativebox}[Approach 3: hand out the blocks top-down --- no stack at all]
Because $p_i < i$, every vertex is processed after its parent in the loop $u = 1..n$. If $u$ already knows where its block starts, its children's blocks follow right after $u$, one after another in increasing index: child $c_1$ starts at $\mathrm{pos}[u]+1$, child $c_2$ at $\mathrm{pos}[u]+1+\mathrm{sz}[c_1]$, and so on. Sizes come first, bottom-up.

\begin{lstlisting}[style=cpp]
#include <bits/stdc++.h>
using namespace std;

// Alternative without any stack: because p_i < i, compute the sizes bottom-up (i = n..2),
// then hand out the blocks top-down (u = 1..n): u's children get consecutive blocks
// starting right after u, in increasing index order.
int main() {
    ios::sync_with_stdio(false);
    cin.tie(nullptr);
    int n, q;
    cin >> n >> q;
    vector<int> p(n + 1, 0);
    vector<vector<int>> ch(n + 1);
    for (int i = 2; i <= n; i++) {
        cin >> p[i];
        ch[p[i]].push_back(i);
    }
    vector<int> sz(n + 1, 1), pos(n + 1, 0), order(n);
    for (int i = n; i >= 2; i--) sz[p[i]] += sz[i];     // children have larger indices
    pos[1] = 0;
    for (int u = 1; u <= n; u++) {                      // pos[u] is known: its parent came earlier
        order[pos[u]] = u;
        int next = pos[u] + 1;                          // first free slot inside u's block
        for (int c : ch[u]) { pos[c] = next; next += sz[c]; }
    }
    while (q--) {
        int u, k;
        cin >> u >> k;
        cout << (k <= sz[u] ? order[pos[u] + k - 1] : -1) << "\n";
    }
}
\end{lstlisting}

\begin{diagrambox}
\begin{tikzpicture}[x=0.95cm, y=1cm]
    \foreach \i in {0,...,8} \node[gcell, minimum size=7mm, font=\scriptsize] (b\i) at (\i,0) {};
    \foreach \i in {0,...,8} \node[font=\tiny, text=darkgray] at (\i,-0.55) {\i};
    \node[gcell, minimum size=7mm, font=\scriptsize\bfseries, fill=currentfill] at (0,0) {1};
    \draw[successgreen, line width=1.2pt] (0.55,0.45) -- (1.45,0.45) node[midway, above, font=\scriptsize] {2 (sz 1)};
    \draw[accentorange, line width=1.2pt] (1.55,0.45) -- (7.45,0.45) node[midway, above, font=\scriptsize] {3 (sz 6)};
    \draw[softpurple, line width=1.2pt] (7.55,0.45) -- (8.45,0.45) node[midway, above, font=\scriptsize] {4 (sz 1)};
    \node[note, anchor=west, text width=5.2cm, align=left] at (9.0,0.1) {$u = 1$: pos$[2] = 1$, pos$[3] = 1+1 = 2$, pos$[4] = 2+6 = 8$. Later $u = 3$ splits $2..7$ into $3 \mid 5,6,8 \mid 7,9$.};
\end{tikzpicture}
\figcap{Approach 3 on the sample: a vertex cuts its own block into the blocks of its children, left to right.}
\end{diagrambox}

\begin{complexitybox}
\textbf{Time:} \bigO{n + q}. \textbf{Space:} \bigO{n}. \textbf{Verdict:} \ac{}-ready; it gives the same \texttt{pos} and \texttt{order} as Approach~2 (identical outputs on the $2\cdot10^5$ path test). It works only because $p_i < i$ gives a parents-first order for free; for a tree given by an edge list, use Approach~2 or the template.
\end{complexitybox}
\end{alternativebox}

\begin{tricksbox}[Approach 4: Implementation Refinements]
\begin{itemize}[leftmargin=*]
    \item \textbf{One clock would not do.} $\tin{}$ from the one-clock convention skips numbers (exits use them too); indexing \texttt{order} needs the entries-only counter \texttt{pos}.
    \item \textbf{Fast I/O.} $2\cdot10^5$ answers: \texttt{"\textbackslash n"}, never \texttt{endl}; \texttt{sync\_with\_stdio(false)}.
    \item \textbf{Reverse push.} Forgetting to push the children in reverse produces the pre-order with children by \emph{decreasing} index --- wrong answers, but only on vertices with $\ge 2$ children.
    \item \textbf{Recursive alternative.} A recursive pre-order is three lines, but at depth $2\cdot10^5$ it crashes on a 2\,MB stack; our test of exactly that case terminated abnormally.
\end{itemize}
\end{tricksbox}

\subsubsection*{5. Trace}

\begin{tracebox}[{Sample: children 1: 2 3 4,\ 3: 5 7,\ 5: 6 8,\ 7: 9}]
\begin{center}
\begin{tikzpicture}[x=0.72cm, y=0.72cm]
    \tikzset{gv/.append style={minimum size=6mm, font=\scriptsize\bfseries}}
    \begin{scope}[g2p3m-v1/.style=gvdone, g2p3m-v2/.style=gvcur, g2p3m-v3/.style=gvseen, g2p3m-v4/.style=gvseen]
        \pic {g2p3m-T};
        \node[font=\scriptsize\bfseries] at (2.4,0.9) {pop 2 (pos 1)};
        \node[font=\scriptsize\ttfamily] at (2.4,-4.5) {stack: 4 3};
    \end{scope}
    \begin{scope}[xshift=4.6cm, g2p3m-v1/.style=gvdone, g2p3m-v2/.style=gvdone, g2p3m-v3/.style=gvdone,
                  g2p3m-v5/.style=gvcur, g2p3m-v4/.style=gvseen, g2p3m-v7/.style=gvseen, g2p3m-v6/.style=gvseen, g2p3m-v8/.style=gvseen]
        \pic {g2p3m-T};
        \node[font=\scriptsize\bfseries] at (2.4,0.9) {pop 5 (pos 3)};
        \node[font=\scriptsize\ttfamily] at (2.4,-4.5) {stack: 4 7 8 6};
    \end{scope}
    \begin{scope}[xshift=9.2cm, g2p3m-v1/.style=gvdone, g2p3m-v2/.style=gvdone, g2p3m-v3/.style=gvdone,
                  g2p3m-v5/.style=gvdone, g2p3m-v6/.style=gvdone, g2p3m-v8/.style=gvdone, g2p3m-v7/.style=gvcur,
                  g2p3m-v4/.style=gvseen, g2p3m-v9/.style=gvseen]
        \pic {g2p3m-T};
        \node[font=\scriptsize\bfseries] at (2.4,0.9) {pop 7 (pos 6)};
        \node[font=\scriptsize\ttfamily] at (2.4,-4.5) {stack: 4 9};
    \end{scope}
\end{tikzpicture}
\end{center}
\begin{center}\small
\begin{tabular}{@{}clll@{}}
\toprule
\textbf{pop} & \textbf{pos} & \textbf{push (reversed children)} & \textbf{stack after} \\
\midrule
1 & 0 & 4, 3, 2 & \texttt{4 3 2} \\
2 & 1 & --- & \texttt{4 3} \\
3 & 2 & 7, 5 & \texttt{4 7 5} \\
5 & 3 & 8, 6 & \texttt{4 7 8 6} \\
6 & 4 & --- & \texttt{4 7 8} \\
8 & 5 & --- & \texttt{4 7} \\
7 & 6 & 9 & \texttt{4 9} \\
9 & 7 & --- & \texttt{4} \\
4 & 8 & --- & (empty) \\
\bottomrule
\end{tabular}
\qquad
\begin{tabular}{@{}cccc@{}}
\toprule
\textbf{query} & \textbf{pos, sz} & \textbf{index} & \textbf{answer} \\
\midrule
$(3,1)$ & 2, 6 & 2 & 3 \\
$(1,5)$ & 0, 9 & 4 & 6 \\
$(3,4)$ & 2, 6 & 5 & 8 \\
$(7,3)$ & 6, 2 & $3 > 2$ & $-1$ \\
$(1,8)$ & 0, 9 & 7 & 9 \\
$(1,9)$ & 0, 9 & 8 & 4 \\
\bottomrule
\end{tabular}
\end{center}
Sizes from the bottom-up loop: $\mathrm{sz} = (9, 1, 6, 1, 3, 1, 2, 1, 1)$ for officers $1..9$. Orange $=$ just popped, light orange $=$ on the stack, blue $=$ already numbered.
\end{tracebox}

\subsubsection*{6. Test suite}

\begin{testbox}
\centering\small
\begin{tabular}{@{}p{6.6cm}lp{4.8cm}@{}}
\toprule
\textbf{Input} & \textbf{Our output} & \textbf{Why} \\
\midrule
\texttt{2 3 / 1 / 1 1 / 1 2 / 2 2} & \texttt{1 2 -1} & smallest tree; $k$ beyond a leaf \\
\texttt{5 5 / 1 2 3 4 / 1 5 / 2 4 / 3 3 / 5 1 / 5 2} & \texttt{5 5 5 5 -1} & a path: every block ends at 5 \\
\texttt{5 4 / 1 1 1 1 / 1 2 / 1 5 / 2 1 / 3 2} & \texttt{2 5 2 -1} & a star: children by index \\
\texttt{6 3 / 1 1 2 2 3 / 1 6 / 2 2 / 3 2} & \texttt{6 4 6} & the subtree of 2 comes before 3 \\
path $p_i = i-1$, $n = q = 2\cdot10^5$ & (same for 2 and 3) & depth $2\cdot10^5$, iterative only \\
\bottomrule
\end{tabular}
\end{testbox}
All outputs come from running Approach~2. Approaches 1--3 reproduce the official sample and agree on $1000$ random trees with $n \le 10$ (paths, stars and random parents, $q \le 10$ random queries each).

\subsubsection*{7. Common mistakes}

\begin{warningbox}[Mistake 1 --- one DFS per query]
Re-running the spreading (even iteratively) for every query is the brute force: \bigO{nq}. The point of the problem is that all $n$ lists are pieces of one list.
\end{warningbox}

\begin{warningbox}[Mistake 2 --- \texttt{tin} instead of \texttt{pos}]
With the one-clock times of the lesson, $\tin{3} = 3$ on the sample, but officer 3 is at position 2 of the pre-order; \texttt{order[tin[u]+k-1]} even runs past the end of the array for large $k$.
\end{warningbox}

\begin{warningbox}[Mistake 3 --- off-by-one in the bound]
$u$ himself is the first officer: the block has $\mathrm{sz}[u]$ elements and the answer exists iff $k \le \mathrm{sz}[u]$, at index $\mathrm{pos}[u] + k - 1$. Query $(7, 3)$ of the sample ($\mathrm{sz}[7] = 2$) catches \texttt{k <= sz[u] + 1}.
\end{warningbox}

\begin{warningbox}[Mistake 4 --- children in the wrong order]
Pushing the children in increasing order onto a stack pops the \emph{largest} first; the sample then prints \texttt{3 9 5 -1 6 2} instead of \texttt{3 6 8 -1 9 4} (our run of that bug). Only the first answer, which never leaves officer 3 himself, survives. Reverse the push, or use the iterator form of the template.
\end{warningbox}

\subsubsection*{8. Takeaways}

\begin{itemize}[leftmargin=*]
    \item A DFS started at $u$ visits exactly the part of the global pre-order that belongs to $u$'s subtree: compute the global order once.
    \item Subtree $=$ block \texttt{order[pos[u] .. pos[u]+sz[u]-1]}: positions count entries only.
    \item With $p_i < i$, sizes need only a reverse loop, and even the positions can be handed out top-down without a stack.
    \item $n = 2\cdot10^5$ and a parent array allow a path: iterative.
\end{itemize}

\begin{variantsbox}
\variant{The $k$-th officer at a given depth}{\varnone}
  {query $(u, d, k)$: the $k$-th officer of depth $d$ that receives $u$'s command.}
  {the depth-$d$ vertices of $u$'s subtree are a contiguous run of the depth-$d$ list sorted by $\mathrm{pos}$.}
  {binary search the first entry $\ge \mathrm{pos}[u]$ in that list and step $k-1$ further; check that it is still $< \mathrm{pos}[u]+\mathrm{sz}[u]$ (the idea of Problem~\ref{ssec:ch2-count-descendants}).}
\variant{Where in the list is officer $v$?}{\varnone}
  {query $(u, v)$: at which step does $v$ receive $u$'s command, or never?}
  {$v$ receives it iff $v$ is in the block of $u$.}
  {answer $\mathrm{pos}[v] - \mathrm{pos}[u] + 1$ if $\mathrm{pos}[u] \le \mathrm{pos}[v] < \mathrm{pos}[u]+\mathrm{sz}[u]$, else ``never''; checked against the brute force on random trees.}
\variant{Command returns to the issuer}{\varin{https://atcoder.jp/contests/abc213/tasks/abc213_d}}
  {every officer reports back to his superior when done; list every visit, including returns.}
  {the Euler walk: a vertex appears on entry and after each child.}
  {print on push and again after each pop (Exercise E3.2).}
\variant{Children by decreasing seniority}{\varnone}
  {each officer passes the command to his subordinates by a given priority instead of index.}
  {same tree, different child order.}
  {sort each \texttt{ch[u]} by the priority before the pre-order; queries unchanged.}
\end{variantsbox}
